\documentclass{article}

% ---------------------------------------------------------
% Professional Formatting
% ---------------------------------------------------------
\usepackage[a4paper, margin=1in]{geometry}
\usepackage{setspace}
\usepackage{parskip}        % paragraph spacing
\setstretch{1.2}            % line spacing

% ---------------------------------------------------------
% Math & Symbols
% ---------------------------------------------------------
\usepackage{amsmath, amssymb, amsthm, mathtools, bm}

% ---------------------------------------------------------
% Figures & Tables
% ---------------------------------------------------------
\usepackage{graphicx}
\usepackage{float}
\usepackage{caption}
\usepackage{subcaption}
\usepackage{booktabs}

% ---------------------------------------------------------
% Code Blocks
% ---------------------------------------------------------
\usepackage{listings}
\usepackage{xcolor}

% ---------------------------------------------------------
% Hyperlinks
% ---------------------------------------------------------
\usepackage{hyperref}
\hypersetup{
    colorlinks=true,
    linkcolor=black,
    urlcolor=blue,
    citecolor=blue
}



% ---------------------------------------------------------
% Title
% ---------------------------------------------------------
\title{
    \vspace{-1cm}
    \textbf{Lab Report}\\
    \vspace{0.2cm}
    Summer 2026
}
\author{
    Soeun Jang \\
    Virginia Tech \\
    \texttt{sophiaj04@vt.edu}
}
\date{\today}

\makeatletter
\renewcommand{\maketitle}{
  \begin{center}
    {\LARGE \bfseries \@title \par}
    \vspace{0.5em}
    \rule{0.8\textwidth}{0.4pt}\par
    \vspace{0.5em}
    {\large \@author \par}
    \vspace{0.5em}
    {\@date}
  \end{center}
}
\makeatother


\pagestyle{plain}


% ---------------------------------------------------------
% Document
% ---------------------------------------------------------
\begin{document}

\maketitle
\newpage   

\tableofcontents
\newpage   

\section*{Introduction}
\addcontentsline{toc}{section}{Introduction}

This report summarizes key concepts from Lab 1–3: low‑rank SVD, Taylor series time stepping, Step‑and‑Truncate methods, Tucker tensors, and low‑rank PDE solvers. Each lab focuses on the methods to express high dimensional data effectively and compute the data.

\begin{itemize}
    \item High dimensional PDEs are too large to compute directly.
    \item Therefore, we compress them using low-rank structures such as SVD or Tucker.
    \item These compressed forms allow efficient evolution of PDE solutions.
\end{itemize}

The workflow across all labs follows the same pattern:

\begin{enumerate}
    \item Compress the object (matrix or tensor).
    \item Apply the PDE operator in low-rank form.
    \item Truncate to prevent rank explosion.
    \item Use Taylor or RK4 for time stepping.
    \item Analyze how the rank changes over time.
    \item Use low-rank structure for gradient computation or extend to 3D with Tucker.
\end{enumerate}

The goal of these labs is to compute high-dimensional PDEs fast while keeping computational cost low. Real PDE solutions mainly contain low-rank structure, meaning they can be represented accurately using compressed forms such as SVD or Tucker.

\section*{Low-Rank SVD}
\addcontentsline{toc}{section}{Low-Rank SVD}

This is the matrices in low rank SVD form and helps to learn how to truncate them efficiently without computing full SVD. Low-Rank helps to reduce computational cost ($O(nr)$ becomes $O(n^2)$) which means data is focused on a few important singular values. This also saves the storage and deletes the noise.

\textbf{Pros:}
\begin{itemize}
    \item Reduces computational cost.
    \item Saves storage.
    \item Removes noise by focusing on dominant singular values.
\end{itemize}

\textbf{Cons:}
\begin{itemize}
    \item Full SVD is expensive.
    \item Hard to choose the correct rank.
\end{itemize}


\section*{Step-and-Truncate}
\addcontentsline{toc}{section}{Step-and-Truncate}

This is the method that while processing time stepping method by taylor series, it helps to prevent the rank getting bigger with the truncation every steps. Since it keeps the rank, its computational cost is consistent so it can save the memories, which means we can simulate for long term.

\textbf{Pros:}
\begin{itemize}
    \item Prevents rank explosion.
    \item Keeps computational cost consistent.
    \item Enables long-term simulation.
\end{itemize}

\textbf{Cons:}
\begin{itemize}
    \item Truncation introduces error.
    \item Frequent truncation increases cost.
\end{itemize}

\section*{Taylor series}
\addcontentsline{toc}{section}{Taylor series}

The n-th order taylor expansion improves the PDE solution. It does contain high accuracy but its computational cost is too expensive. This is accurate method and simple to work, but repeated calculation of truncation makes the speed slow (especially in tucker structure).

\textbf{Pros:}
\begin{itemize}
    \item High accuracy.
    \item Simple to implement.
\end{itemize}

\textbf{Cons:}
\begin{itemize}
    \item Computationally expensive.
    \item Repeated truncation slows performance.
\end{itemize}


\section*{Strategy A vs Strategy B}
\addcontentsline{toc}{section}{Strategy A vs Strategy B}

A represents truncating each terms and B represents truncating only once at the end. A is more accurate to keep rank low but B is faster to compute since it does not have to compute until the end.

\textbf{Pros of Strategy A:}
\begin{itemize}
    \item More stable.
    \item Keeps rank low.
\end{itemize}

\textbf{Cons of Strategy A:}
\begin{itemize}
    \item Slower due to frequent truncation.
\end{itemize}

\textbf{Pros of Strategy B:}
\begin{itemize}
    \item Faster computation.
\end{itemize}

\textbf{Cons of Strategy B:}
\begin{itemize}
    \item Risk of temporary rank growth.
\end{itemize}
\section*{Sensitivity matrix}
\addcontentsline{toc}{section}{Sensitivity matrix}

Computing sensitivity matrix indicates the way how the PDE solution changes by parameters. Since the sensitivity matrix is big, it is necessary to contain low rank structure. However, low-rank structure is needed and Dense structure only stores $O(n^2)$, cannot compute but it can compute its gradient.

\textbf{Pros:}
\begin{itemize}
    \item Shows how solution changes with parameters.
    \item Low-rank structure makes computation feasible.
\end{itemize}

\textbf{Cons:}
\begin{itemize}
    \item Dense sensitivity matrix is too large to store.
\end{itemize}
\section*{Cheap gradient}
\addcontentsline{toc}{section}{Cheap gradient}

Computing gradient using only low-rank factors without full SVD matrices, only with Tucker and SVD factors to reduce the computational cose but gradient can be inaccurate when rank is too small and truncational error might affect the gradient.

\textbf{Pros:}
\begin{itemize}
    \item Fast and memory-efficient.
    \item Works well in optimization loops.
\end{itemize}

\textbf{Cons:}
\begin{itemize}
    \item Inaccurate when rank is too small.
    \item Truncation error may affect gradient.
\end{itemize}

\section*{Tensor Operations}
\addcontentsline{toc}{section}{Tensor Operations}

Since 3D tensor contains too big size, we saves this tensor as small core and basis by compressing efficiently. It saves space to store and computational costs, also we can observe its mode structures. However, it is complicated. In step-and-truncate in tucker, we keep low-rank in compressed form since dense tensor is too big when we solve 3D PDE. This costs a lot by truncating especially in taylor method but it keeps rank stable and accuracy with reducing costs for 3D PDE.

\textbf{Pros:}
\begin{itemize}
    \item Saves storage.
    \item Reduces computational cost.
    \item Reveals mode structure.
\end{itemize}

\textbf{Cons:}
\begin{itemize}
    \item Tucker operations are complex.
    \item Truncation is expensive.
\end{itemize}


\section*{Conclusion}
\addcontentsline{toc}{section}{Conclusion}

Lab 1-3 allowed us to see directly why the low-rank method is important and how effective it is in real PDE calculations. Beginning with the matrix SVD, solving 2D PDE and expanding to 3D PDE with tucker tensor. Based on the results, many PDE solutions and sensitivity matrices actually contain low-rank structure, so it was impossible to calculate accurately enough with only a small rank without computing the full SVD and tensors. Additionally, it keeps the rank low with the proper use of truncation and maintain the accuracy. In conclusion, low-rank SVD represents the tool as solving high dimensional problem fast and efficient.

\end{document}
